% © 2026 Ing. David Jaroš — CC BY-NC-ND 4.0
\documentclass[12pt,a4paper]{article}
\usepackage[a4paper,margin=2.5cm]{geometry}
\usepackage{amsmath,amssymb,amsthm,mathtools}
\usepackage[most]{tcolorbox}
\newtheorem{definition}{Definition}[section]
\newtheorem{theorem}[definition]{Theorem}
\newtheorem{corollary}[definition]{Corollary}
\theoremstyle{remark}\newtheorem{remark}[definition]{Remark}
\newcommand{\B}{\mathbb B}
\title{GAP-10L and GAP-10$\psi$: Symmetry Propagation and Metric Stability}
\author{Ing. David Jaro\v{s}}
\date{16 July 2026}
\begin{document}
\maketitle

\begin{abstract}
This note identifies exact sufficient conditions for preservation of the UBT
Lorentz slice and for independence of the classical metric from imaginary
time.  The Lorentz slice is the fixed set of a natural anti-linear involution
$\mathcal JX=-\overline{X^\sharp}$.  Any well-posed evolution whose equations,
sources and data are $\mathcal J$-equivariant preserves that fixed set by
uniqueness.  Independently, a $\psi$-flow of the tetrad tangent to a local
Lorentz gauge orbit leaves the anticommutator metric exactly unchanged.  A
$\psi$-translation-invariant well-posed evolution also preserves strictly
$\psi$-independent data.  These theorems close symmetry-propagation subgaps;
verification that the final canonical UBT action satisfies their hypotheses
remains open.
\end{abstract}

The fixed-set propagation argument used below is the standard consequence of
equivariance plus uniqueness in a well-posed dynamical system; compare the
general symmetry framework for differential equations in
Olver~\cite{olver1993}.  The UBT-specific tasks are to identify the intrinsic
involution and to verify that the eventual canonical equations satisfy the
required hypotheses.

\section{The real Lorentz slice as a fixed set}
Write a biquaternion as $X=a_0\mathbf1+a_k\mathbf e_k$ with $a_\alpha\in\mathbb C$
and quaternionic conjugation
$X^\sharp=a_0\mathbf1-a_k\mathbf e_k$.  Define
\begin{equation}
 \boxed{\mathcal JX:=-\overline{X^\sharp}.}
 \label{eq:J}
\end{equation}

\begin{theorem}[Intrinsic characterization of the Lorentz slice]
$\mathcal J$ is an anti-linear involution and
\begin{equation}
 \operatorname{Fix}(\mathcal J)
 =\{i x^0\mathbf1+x^k\mathbf e_k\mid x^a\in\mathbb R\}
 =W_L.
\end{equation}
\end{theorem}

\begin{proof}
For $X=a_0\mathbf1+a_k\mathbf e_k$,
$\mathcal JX=-\bar a_0\mathbf1+\bar a_k\mathbf e_k$; applying $\mathcal J$
twice returns $X$.  The fixed-point equations are
$a_0=-\bar a_0$ and $a_k=\bar a_k$, hence $a_0$ is purely imaginary and the
three vector coefficients are real.
\end{proof}

\section{Equivariance implies dynamical slice preservation}
Let $U$ collect the fields needed for a Cauchy formulation, including $\Theta$
and, when necessary, its canonical momentum or first time derivative.

\begin{theorem}[Fixed-set propagation by uniqueness]
Assume a Cauchy problem
\begin{equation}
 \partial_t U=\mathcal F(U;S)
 \label{eq:evolution}
\end{equation}
is locally well posed and unique in a function class invariant under
$\mathcal J$.  Suppose
\begin{equation}
 \mathcal F(\mathcal JU;\mathcal JS)
 =\mathcal J\mathcal F(U;S),
 \label{eq:equivariance}
\end{equation}
and the source and initial data are fixed by $\mathcal J$.  Then the unique
solution remains fixed by $\mathcal J$ throughout its interval of existence.
\end{theorem}

\begin{proof}
If $U$ solves Eq.~\eqref{eq:evolution}, equivariance implies that
$\mathcal JU$ is another solution.  The two solutions have identical initial
data and sources.  Uniqueness therefore gives $\mathcal JU=U$.
\end{proof}

\begin{corollary}[Lorentz-slice preservation]
If the complete UBT equations induce a $\mathcal J$-equivariant unique
evolution for the tetrad variables $E_\mu=D_\mu\Theta/\sqrt{\mathcal N_0}$,
then initial data $E_\mu\in W_L$ and $\mathcal J$-real sources imply
$E_\mu(t)\in W_L$ for all times for which the solution exists.
\end{corollary}

\section{Imaginary-time metric stability}
Let $e_\mu{}^a(x,\psi)$ be the real tetrad coefficients and
$g_{\mu\nu}=e_\mu{}^ae_\nu{}^b\eta_{ab}$.

\begin{theorem}[$\psi$-vertical Lorentz flow leaves the metric invariant]
If
\begin{equation}
 \partial_\psi e_\mu{}^a
 =\Lambda_\psi{}^a{}_b e_\mu{}^b,
 \qquad \Lambda_{\psi ab}=-\Lambda_{\psi ba},
 \label{eq:vertical-flow}
\end{equation}
then
\begin{equation}
 \boxed{\partial_\psi g_{\mu\nu}=0.}
\end{equation}
Thus strict $\psi$-independence of the tetrad is sufficient but not necessary:
its $\psi$ dependence may be pure local Lorentz gauge.
\end{theorem}

\begin{proof}
Differentiate $g_{\mu\nu}$.  The two terms are
$\Lambda_{\psi ab}e_\mu{}^be_\nu{}^a$ and
$\Lambda_{\psi ba}e_\mu{}^be_\nu{}^a$, which cancel by antisymmetry.
\end{proof}

\begin{theorem}[$\psi$-translation symmetry propagation]
Suppose the complete Cauchy/boundary problem is invariant under
$\psi\mapsto\psi+s$ and has a unique solution.  If all initial and boundary
data and sources are $\psi$ independent, then the solution is $\psi$
independent.
\end{theorem}

\begin{proof}
For every $s$, the translated field $U_s(x,\psi)=U(x,\psi+s)$ is a solution
with the same data.  Uniqueness gives $U_s=U$ for all $s$, hence
$\partial_\psi U=0$.
\end{proof}

\section{Status}
\begin{tcolorbox}[colback=green!6!white,colframe=green!55!black,title=Defensible status]
\textbf{GAP-10L-SYM: CLOSED CONDITIONALLY [L1].}
The Lorentz slice is the fixed set of the intrinsic involution
$\mathcal J=-\overline{(\cdot)^\sharp}$, and every unique
$\mathcal J$-equivariant evolution with fixed data preserves it.

\medskip
\textbf{GAP-10L-DYN: NARROWED.}
It remains to prove that the finalized UBT action, its paired connections,
sources and chosen well-posed formulation satisfy the equivariance and
uniqueness hypotheses.

\medskip
\textbf{GAP-10$\psi$-KIN: CLOSED [L1].}
A $\psi$-flow tangent to the local Lorentz gauge orbit leaves the metric
exactly invariant.

\medskip
\textbf{GAP-10$\psi$-SYM: CLOSED CONDITIONALLY [L1].}
A unique $\psi$-translation-invariant evolution preserves $\psi$-independent
data.

\medskip
\textbf{GAP-10$\psi$: NARROWED.}
The canonical action must still show that the physical classical branch obeys
one of these sufficient mechanisms and that no unstable non-gauge
$\psi$ excitations contaminate the observed metric.
\end{tcolorbox}

\begin{thebibliography}{9}
\bibitem{olver1993}
P.~J.~Olver,
\emph{Applications of Lie Groups to Differential Equations}, 2nd ed.,
Springer, New York (1993).
\end{thebibliography}
\end{document}
